\section{Geometric Deep Learning}

\begin{definition}[Symmetry]
    Transformation that leaves an object invariant.
\end{definition}

\begin{definition}[Group]
    Set \(\mathfrak{G}\) with composition \(\circ: \G \times \G \to \G\), satisfying:
    \begin{itemize}
        \item Associativity: \(\forall \gg, \hg, \fg \in \G: (\gg\hg)\fg = \gg(\hg\fg)\).
        \item Identity: \(\exists\) unique \(\eg \in \G \forall \gg: \eg\gg = \gg\eg = \gg\).
        \item Inverse: \(\forall \gg \exists\) unique \(\gg^{-1}: \gg\gg^{-1} = \gg^{-1}\gg = \eg\)
        \item Closure: \(\forall \gg, \hg \in \G: \gg\hg \in \G\).
    \end{itemize}
\end{definition}

\begin{definition}[Signal Space]
    \(\DX(\Omega, \mathcal{C}) := \{x: \Omega \to \mathcal C\}\) space of \(\mathcal C\)-valued signals on set \(\Omega\). Induces linear structure.
\end{definition}

\begin{definition}[Action]
    Mapping \((\gg, u) \mapsto \gg.u\), w/ \(\gg.(\hg.u) = (\gg\hg).u\), mapping group elements to transformations on set \(\Omega \ni u\). Implicitly defines action of \(\G\) onto signals: \((\gg.x)(u) := x(\gg^{-1}.u)\) with \((\gg(\hg.x))(u) = ((\gg\hg).x)(u)\).
\end{definition}

\begin{definition}[Representation]
    Represent linear actions of \(\G\) via \(\rho: \G \to \R^{n \times n}\), assing invertible matrix \(\rho(\gg)\) and satisfying \(\rho(\gg\hg) = \rho(\gg) \rho(\hg)\). Again, stated \(\rho(\gg)x(u) = x(\gg^{-1} u)\) and \((\rho(\gg)(\rho(\hg)x))(u) = (\rho(\gg\hg)x)(u)\).
\end{definition}

\begin{definition}[\(\G\)-invariance]
    \(\forall \gg \in \G, x\in\DX(\Omega): f(\rho(\gg)x) = f(x)\).
\end{definition}

\begin{definition}[\(\G\)-equivariance]
    \(\forall \gg \in \G, x\in\DX(\Omega): f(\rho(\gg)x) = \rho(\gg)f(x)\).
\end{definition}

\begin{colored}
    \begin{itemize}
        \item \(f\) non-invaraint \(\to\) \(f^*\) invariant: \(S_\G f(x) = \frac{1}{|\G|}\sum_{g \in \G} f(\rho(\gg)(x))\)
    \end{itemize}
\end{colored}

\begin{tabularx}{\linewidth}{XXX}
\textbf{Architecture} & \textbf{Domain} & \textbf{Symmetry} $\mathfrak{G}$ \\
\hline
CNN & Grid & Translation \\ \hline
Spherical CNN & Sphere / SO(3) & Rotation SO(3) \\ \hline
Deep Sets & Set & Permutation $\Sigma_n$ \\ \hline
GNN & Graph & Permutation $\Sigma_n$ \\ \hline
Transformer & Complete Graph & Permutation $\Sigma_n$ \\ \hline
LSTM & 1D Grid & Time warping \\ \hline
\end{tabularx}

\begin{definition}[5-Gs]
    Grids, Groups, Graphs, Geodesics, Gauges.
\end{definition}

\subsection{Set Neural Networks}
Set domain \(V\) with \(\X = (\x_1, \ldots, \x_n)^\top \in \R^{n\times d}\), \(\G = \Sigma_n\), \(\rho(\gg) = \PM\).

\begin{definition}[Permuation invariance]
    \(\forall P \in \Sigma_n: f(\PM\X) = f(\X) = \phi\left(\sum_i \psi(\x_i)\right)\)
\end{definition}

\begin{definition}[Permutation equivariance]
    \(\forall P \in \Sigma_n: f(\PM\X) = \PM f(\X)\), could be implemented as \(f_i(\X) = f(\X, \x_i) = \phi(\x_i, \sum_{j} \psi(\x_j))\).
\end{definition}

\subsection{Graph Neural Networks}
Graph \(G = (V, E, \A)\) w/ \(a_{ij} > 0\) if \((i, j) \in E\), else \(0\).
Rest assumes \textbf{undirected weighted graph}.

\begin{definition}[Canonical Graph]
    \(G_o = (V, E_0)\) with \(E_o = \{(i, j) \in E \mid i < j\}\).
\end{definition}

\begin{definition}[Incidence matrix]
    \(\B \in \R^{|E_o| \times |V|}\) defined for edge \(e = (i, j) \in E_o\) by \(b_{e, v} = 1\) if \(v = i\), \(-1\) if \(v = j\), else \(0\).
\end{definition}

\begin{definition}[Degree Matrix]
    \(\DM = \text{diag}(d_i)\) with \(d_i = \sum_j a_{ij}\).
\end{definition}

\begin{definition}[Edge Weight Matrix]
    \(\A_e = \text{diag}(a_e)\) with \(a_e = a_{ij}\).
\end{definition}

\begin{lemma}
    \begin{itemize}
        \item \(\B^\top \A_e \B = \DM - \A\)
        \item \(G\) unweighted \(\iff\) \(\A_e\) identity
        \item \(G\) undirected \(\iff\) \(\A = \A^\top\).
    \end{itemize}
\end{lemma}

\begin{definition}[Invariance for UWG]
    \(f(\PM\X, \PM\A\PM^\top) = f(\X, \A)\)
\end{definition}

\begin{definition}[Equivariance for UWG]
    \(f(\PM\X, \PM\A\PM^\top) = \PM f(\X, \A)\)
\end{definition}

\begin{definition}[Neighborhood]
    \(\Normal_i := \{j : (i, j) \in E\}\), \(\X_{(\Normal_i)} = \multiset{\x_i: j \in \Normal_i}\)
\end{definition}

\begin{definition}[Layer]
    \(\phi(\x_i, \X_{(\Normal_i)})\) perm. inv. in \(\X_{(\Normal_i)}\). \(f(\X, \A) = [\phi(\x_i, \X_{(\Normal_i)})]_i\)
\end{definition}

\begin{definition}[Equivariant Layer]
    \(\phi(\x_i, \X_{(\Normal_i)}) = \phi\left(\x_i, \sum_{j \in \Normal_i} \psi(\x_j)\right)\)
\end{definition}

\begin{definition}[Conv. GNNs]
    \(z_i = \phi\left(\x_i, \sum_j c_{ij}(\A) \psi(\x_j)\right)\) w/ \(\CM_{\A} = (c_{ij}(\A))_{ij}\) s.t.
    \(\CM_{\PM \A \PM^\top} = \PM \CM_{\A} \PM^\top\) self-equivariant.
\end{definition}

\begin{definition}[Attentional GNNs]
    \(z_i = \phi\left(\x_i, \sum_{j \in \Normal_i} a(\x_i, \x_j) \psi(\x_j)\right)\).
\end{definition}

\begin{definition}[Message-Passing GNNs]
    \(z_i = \phi\left(\x_i, \sum_{j \in \Normal_i} \psi(\x_i, \x_j)\right)\).
\end{definition}

\begin{definition}[Weisfeiler-Lehman Test]
    Test if two graphs are non-isomorphic by iteratively aggregating node features from neighbors and hashing them to new features. If after some iterations the multisets of node features differ, the graphs are non-isomorphic. NP-hard to decide in general.
\end{definition}

\begin{theorem}[Aggregators]
    In order to discriminate multisets of size \(n\), whose underlying set is \(\R\), at least \(n\) aggregators are needed.
\end{theorem}

\subsection{Spectral Graph Theory (\(n := |V|, m := |E_o|\))}

\begin{definition}[Discrete Derivative]
    Along edge \((i, j)\): \(\partial f_{ij} = \sqrt{a_{ij}}(f_j - f_i)\).
\end{definition}

\begin{definition}[Divergence]
    At vertex \(i\): \([\text{div}(g)]_i = \sum_{j: (i, j) \in E} \sqrt{a_{ij}} g_{ij}\) for \(g: E_0 \to \R\), with convention that \(g_{ij} = - g_{ji}\).
\end{definition}

\begin{definition}[Representations]
    \(\text{grad}(f) = (\partial_{ij} f)_{(i,j) \in E_o}\), \(\text{div}(g) = ([\text{div}(g)]_i)_{i \in V}\). Then \(\grad := \A_e^{1/2} \B \in \R^{m \times n} \implies \grad f = \text{grad}(f)\) and \(\text{div}(g) = \grad^\top g\).
\end{definition}

\begin{definition}[Graph laplacian]
    \(\LM := \grad^\top \grad = \B^\top \A_e \B = \DM - \A\)
\end{definition}

\begin{theorem}
    \([\LM \vb f]_i = [\B^\top \A_e \B \vb f] = \sum_{j \in \Normal_i} a_{ij}(f_i - f_j)\) and \(\vb f^\top \LM \vb f = \norm{\grad \vb f}^2 = \sum_{ij} a_{ij} (f_i - f_j)^2\). \(\LM\) is symmetric positive semi-definite with eigenvalues \(0 = \lambda_1 \leq \lambda_2 \leq \ldots \leq \lambda_n\).
    Hence \(\LM = \U \Lambda \U^\top\) and columns of \(\U\) form orthonormal basis of eigenvectors in fourier domain.
\end{theorem}

\begin{definition}[Spectral Graph Convolution]
    \(\h *_G \x := \U \text{diag}(\hat \h) \U^\top \x\), but requires eigenvalues of decomposition. Hence just want \(\h *_G \x = \U h(\Lambda) \U^\top \x = h(\LM)\x\).

    For efficiency, use polynomial filters \(h(\lambda) = \sum_{k=0}^K \theta_k \lambda^k\) (Chebyshev polynomials) or \(h(\lambda) = \theta (1 - \lambda)\) (first-order approx.).

    In practice, learn \(C_{in} C_{out}\) poly kernels and use \(\x_{c'} = \sum_{c=1}^{C_{in}} p_{c', c}(\LM)\x_c + b_{c'}\vb{1}\) with \(C_{in}C_{out}(K+1)+C_{out}\) params.
\end{definition}

\begin{definition}[Norm. Graph Laplacian]
    \(\bar\LM := \DM^{-1/2} \LM \DM^{-1/2} = \I - \DM^{-1/2}\A \DM^{-1/2}\).
\end{definition}

\begin{definition}[Data Representation]
    Data matrix $\mathbf{X} \in \mathbb{R}^{d \times n}$ where $n$ = nodes, $d$ = channels.
    \textbf{Propagation}: $\mathbf{X}^{l+1} = \sigma(\mathbf{W}^l \mathbf{X}^l \mathbf{Q})$ where $\mathbf{Q}$ couples node activities.
\end{definition}

\begin{definition}[Normalized Graph Laplacian]
    Adjacency $\mathbf{A}$, degree matrix $\mathbf{D}$. 
    $\bar{\mathbf{L}} = \mathbf{I} - \mathbf{D}^{-1/2} \mathbf{A} \mathbf{D}^{-1/2}$.
    \textbf{Properties}: $\lambda(\bar{\mathbf{L}}) \in [0, 2]$. $\lambda_{max} = 2$ can cause instabilities.
\end{definition}

\begin{definition}[GCN Coupling Matrix $\mathbf{Q}$]
    To fix instabilities, use \textbf{Renormalization Trick}: 
    $\tilde{\mathbf{A}} = \mathbf{A} + \mathbf{I}_n$ (self-loops) and $\tilde{\mathbf{D}}$ is the degree matrix of $\tilde{\mathbf{A}}$.
    $\mathbf{Q} = \tilde{\mathbf{D}}^{-1/2} \tilde{\mathbf{A}} \tilde{\mathbf{D}}^{-1/2}$.
    \textbf{Result}: $\lambda_{max}(\mathbf{Q}) = 1$, preventing exploding gradients during stacking.
\end{definition}

\begin{theorem}[Spectral Filtering]
    Using $\mathbf{Q} = \bar{\mathbf{A}}$ (normalized adjacency) is equivalent to a first-order spectral filter $h(\bar{\mathbf{L}}) = \mathbf{I} - \bar{\mathbf{L}}$.
    Stacking $L$ layers allows learning higher-order polynomial filters.
\end{theorem}
